\chapter*{Kết luận}
\addcontentsline{toc}{chapter}{Kết luận}

Luận văn này đã tập trung giải quyết bài toán kiểm tra lỗi màn hình điện thoại trong môi trường sản xuất công nghiệp, đối mặt với thách thức kép về sự thiếu hụt dữ liệu gán nhãn và yêu cầu hệ thống phải có khả năng học liên tục để thích nghi với các loại lỗi mới. Thông qua quá trình nghiên cứu lý thuyết và thực nghiệm, luận văn đã đạt được những đóng góp khoa học và thực tiễn quan trọng. Trước hết, nghiên cứu đã xây dựng thành công quy trình sinh dữ liệu giả lập dựa trên mô hình hóa hình học và nhiễu thống kê, giải quyết triệt để vấn đề khan hiếm dữ liệu đầu vào. Tiếp đó, phương pháp trích chọn đặc trưng lai ghép giữa bản đồ Entropy cục bộ và mạng Autoencoder đã được chứng minh hiệu quả trong việc nâng cao khả năng phát hiện các lỗi có độ tương phản thấp, vốn là điểm yếu của các phương pháp xử lý ảnh truyền thống.

Kết quả nổi bật nhất của đề tài là việc ứng dụng thành công mạng Growing Neural Gas (GNG) cho bài toán học liên tục. Bằng cách tận dụng cơ chế học cạnh tranh và sinh trưởng cấu trúc tô-pô động, hệ thống đã khắc phục hiệu quả hiện tượng "quên thảm họa". Các thực nghiệm qua ba kịch bản tác vụ tuần tự cho thấy mô hình có khả năng tự động mở rộng không gian tri thức để nhận diện các lỗi mới như hình chữ nhật hay hình tam giác, đồng thời vẫn bảo toàn độ chính xác trung bình trên 96\% đối với các lỗi cũ đã học. Bên cạnh đó, tính khả thi của giải pháp cũng được khẳng định thông qua ứng dụng phần mềm demo hoạt động thời gian thực với tốc độ xử lý 45ms/ảnh, hỗ trợ trực quan cho công tác giám sát chất lượng.

Tuy nhiên, nghiên cứu vẫn tồn tại một số hạn chế nhất định cần được nhìn nhận khách quan. Do sự phụ thuộc chủ yếu vào dữ liệu giả lập, mô hình có thể chưa bao quát hết được các nhiễu vật lý phức tạp trong thực tế như vết vân tay, bụi bẩn hay sự thay đổi gắt gao của ánh sáng môi trường. Ngoài ra, chi phí tính toán của thuật toán GNG hiện tăng tuyến tính theo số lượng nơ-ron, điều này có thể gây áp lực lên tài nguyên hệ thống nếu số lượng loại lỗi cần phân biệt tăng lên quá lớn. Các tham số vận hành của mạng hiện tại cũng được thiết lập cố định, chưa thực sự linh hoạt để tự điều chỉnh theo mức độ biến động của dữ liệu đầu vào.

Từ những phân tích trên, các hướng nghiên cứu tiếp theo được đề xuất nhằm hoàn thiện hệ thống bao gồm việc kết hợp cơ chế học bán giám sát để tận dụng dữ liệu thực tế từ dây chuyền, nâng cao độ tin cậy của mô hình. Về mặt thuật toán, nghiên cứu sẽ hướng tới việc tối ưu hóa tốc độ tìm kiếm của GNG bằng các cấu trúc dữ liệu cây hoặc phát triển biến thể GNG phân cấp để quản lý các cụm lỗi hiệu quả hơn. Cuối cùng, mục tiêu dài hạn là tối ưu hóa mô hình để triển khai trên các thiết bị tính toán biên (Edge AI) và camera thông minh, góp phần hiện đại hóa quy trình kiểm soát chất lượng trong các nhà máy thông minh.